\chapter{Практика}
\label{ch:intro}

\section{timing error detection}

\section{Описание}

Этот блок - одна из составных частей схемы Гарднера. Ее задача заключается в том,
чтобы рассчитать некоторый коэффициент, который указывает на ошибку или сдвиг,
который необходимо сделать, чтобы попасть в нужный семпл. В идеале, если изобразить
эти ошибки на графике в зависимости от сдвига в семплах (от 0 до 10), то должны получить
кривую, напоминающую синус (позже проверим)

\subsection{Реализация}

\begin{lstlisting}
std::vector<int16_t> synchronizer::gardner(const std::vector<std::complex<double>> &samples, int const L)
{
    /*define start parameters*/
    double Kp = 4;
    double BnTs = 0.000005;

    double zeta = std::sqrt(2) / 2;
    double theta = (BnTs / 10) / (zeta + 1 / (4 * zeta));
    double den = (1 + 2 * zeta * theta + theta * theta) * Kp;

    double K1 = -4 * zeta * theta / den;
    double K2 = -4 * theta * theta / den;

    double p1 = 0;
    double p2 = 0;

    int offset = L / 2;
    double e = 0;

    int n = 0;

    std::vector<int16_t> offset_list;

    int max_i = samples.size() % L == 0 ? samples.size() / L : samples.size() / L + 1;

    for (int i = 0; i < max_i; ++i)
    {
        n = offset + L * i;
        /*calculate e - timing error*/
        e = (samples[n + L].real() - samples[n].real()) * samples[n + L / 2].real() +
            (samples[n + L].imag() - samples[n].imag()) * samples[n + L / 2].imag();

    }
}
\end{lstlisting}

На Python визуализируем e, чтобы убедиться, что ted вычисляется верно

\begin{figure}[H]
    \centering
    \includegraphics[width=1.0\textwidth]{real_avge.png}
    \caption{TED}
\end{figure}

График напоминает синус, что означает правильность работы TED


\endinput
